% TOP_PROBLEM: {"id": 30006762, "title": "Randomized Integer Word-RAM 3SUM Conjecture", "category_id": 15, "category": {"id": 15, "name": "computer_science", "display_name": "Computer Science", "description": "Computational complexity, algorithms, and theoretical CS.", "slug": "computer-science", "order_index": 15, "created_at": "2026-07-31T15:26:25.671Z"}, "status": "open"}
% FIRST_POSTED: 2026-09-27
% !TeX program = lualatex
% Compile twice: lualatex open_problems_500.tex
% All references and prose are embedded; no BibTeX database or external inputs.
\documentclass[11pt,a4paper]{article}
\usepackage[margin=24mm,headheight=14pt]{geometry}
\usepackage{fontspec}
\setmainfont{Latin Modern Roman}
\setsansfont{Latin Modern Sans}
\setmonofont{Latin Modern Mono}
\usepackage{amsmath,amssymb,mathtools}
\usepackage{unicode-math}
\setmathfont{Latin Modern Math}
\usepackage{microtype}
\usepackage{enumitem,xurl,needspace}
\usepackage[dvipsnames]{xcolor}
\usepackage{hyperref}
\hypersetup{unicode=true,colorlinks=true,linkcolor=MidnightBlue,urlcolor=MidnightBlue,
 pdftitle={500 Mathematical Problem Statements},pdfauthor={Source-annotated compilation},bookmarksnumbered=true}
\usepackage{bookmark}
\usepackage{tocloft}
\setlength{\cftsecnumwidth}{3em}
\cftsetpnumwidth{2.5em}
\cftsetrmarg{4em}
\usepackage{titlesec}
\titleformat{\section}{\Large\bfseries\raggedright}{\thesection}{0.7em}{}
\usepackage{fancyhdr}
\pagestyle{fancy}\fancyhf{}
\fancyhead[L]{\small\sffamily Mathematical problem statements}
\fancyhead[R]{\small Source edition 22}
\fancyfoot[C]{\thepage}
\setlength{\parindent}{0pt}
\setlength{\parskip}{5pt plus 1pt}
\setlength{\emergencystretch}{3em}
\setcounter{tocdepth}{1}
\setcounter{secnumdepth}{1}
\setlist[itemize]{leftmargin=3em,itemsep=5pt,topsep=4pt,parsep=0pt}
\allowdisplaybreaks
\newcommand{\N}{\mathbb{N}}
\newcommand{\Z}{\mathbb{Z}}
\newcommand{\Q}{\mathbb{Q}}
\newcommand{\R}{\mathbb{R}}
\newcommand{\C}{\mathbb{C}}
\newcommand{\F}{\mathbb{F}}
\newcommand{\A}{\mathbb{A}}
\newcommand{\PP}{\mathbb P}
\newcommand{\E}{\mathbb{E}}
\newcommand{\eps}{\varepsilon}
\newcommand{\dd}{\,\mathrm d}
\newcommand{\sref}[2]{\hyperlink{source:#1:#2}{[S#2]}}
\newcommand{\eref}[2]{\hyperlink{extra:#1:#2}{[E#2]}}
% Turn the next switch off for a shorter reading copy. The quotations preserve
% every exactTarget field, including qualifications omitted from a short summary.
\newif\ifshowcatalogue
\showcataloguetrue
\newcommand{\cataloguescope}[1]{\ifshowcatalogue
 \paragraph{Exact scope in the supplied catalogue.}
 \begin{quote}\small #1\end{quote}\fi}
% Compile twice with: lualatex 30006762.tex
\numberwithin{equation}{section}
\hypersetup{pdftitle={Research attempt -- problem 30006762},pdfauthor={Research notebook; original compilation retained}}
\fancyhead[L]{\small\sffamily Research notebook}
\fancyhead[R]{\small\sffamily Problem 30006762 | 27 September 2026}
\begin{document}
\setcounter{section}{0}
% ================================================================
% problemId: 30006762
% Canonical problem ID: 30006762
% exactTarget SHA-256: a94892f0e869f5f358c6b88ab0dea9b2a02ca0eb1ffb31b7751d6173d1a750bb
\clearpage
\section*{\texorpdfstring{Randomized Integer Word-RAM 3SUM Conjecture}{Randomized Integer Word-RAM 3SUM Conjecture}}\label{30006762}
\subsection{Definitions and mathematical statement}
For a set $S\subseteq[-n^3,n^3]\cap\Z$ of $n$ distinct integers, define
$\mathrm{3SUM}(S)=1$ precisely when distinct $a,b,c\in S$ satisfy $a+b+c=0$. Work on a uniform integer word-RAM with $\Theta(\log n)$-bit words and the operations specified by the source. A Las Vegas algorithm always answers correctly and terminates almost surely; its cost is expected running time over its internal random bits. The conjecture excludes, for every fixed $\eps>0$, an algorithm satisfying
\[
 \sup_{|S|=n}\E[T(S)]=O(n^{2-\eps}).
\]
There is no real-arithmetic oracle, free preprocessing, or input-specific advice. Logarithmic improvements to $n^2$ are not excluded by this exponent statement.
\par\smallskip\textit{Formulation and concept references:} \sref{30006762}{1}.
\cataloguescope{For every epsilon>0, there is no uniform Las Vegas randomized integer word-RAM algorithm that, on every set S of n distinct integers in [-n\textasciicircum{}3,n\textasciicircum{}3], decides whether three distinct elements sum to zero in worst-case expected O(n\textasciicircum{}(2-epsilon)) time. Use the standard logarithmic-word model of Abboud–Vassilevska Williams Conjecture 1, with enough Theta(log n) bits for the input universe, unit-cost word arithmetic, comparisons, Boolean operations, shifts and indirect memory access; no real-number oracle or free input-dependent preprocessing is permitted. Las Vegas means an always-correct answer with almost-sure termination, expectation over internal randomness, and a bound uniform over all valid inputs. Constants may depend on the algorithm and epsilon. This is not a literal quadratic lower bound.}
\subsection{Short English statement}
Must exact three-sum testing on bounded integers resist every always-correct randomized algorithm whose worst-case expected exponent is strictly below two?
% BEGIN RESEARCH ADDITION ATTEMPT-161
\subsection{Research attempt}

\paragraph{Current frontier.}
The formulation source treats the exponent form of integer 3SUM hardness,
not an exact quadratic lower bound \sref{30006762}{1}, Conjecture 1. The
logarithmic savings of Gr\o nlund--Pettie do not refute this selected
statement \eref{30006762}{1}. A particularly relevant recent result is the
$O(n^{3/2}\log n)$ query algorithm of Kasliwal--Polak--Sharma after
quadratic preprocessing of three fixed universes
\eref{30006762}{2}. Its preprocessing is forbidden as a free resource here.
The attempt below does not use it. Instead it isolates a streaming
witness-reporting statement whose truth would give a genuinely
subquadratic algorithm in the exact model of the catalogue.

\paragraph{Attack route.}
Dense convolution is fast but the input universe has size $\Theta(n^3)$.
Modular hashing reduces that universe while introducing false triples.
The usual obstruction is that a convolution coefficient reports a count,
not an efficiently accessible list of witnesses. A second obstruction is
that even a true instance may have quadratically many valid witnesses,
although a decision algorithm needs only one. The proposed attack joins a
random modulus with \emph{prefix-efficient} witness enumeration. Both
obstructions are quantified rather than hidden in a subroutine call.

\paragraph{Attempt: the missing algorithmic lemma.}
Fix the following precise hypothesis, denoted $\mathsf{StreamConv}$.
Given a prime $p$ and a Boolean array $A$ indexed by $\Z/p\Z$, there is a
uniform algorithm in the logarithmic-word RAM model that lists, without
repetition, all ordered triples
\begin{equation}\label{30006762:residue-triples}
 (r,s,t),\qquad r+s+t=0\pmod p,\qquad A(r)A(s)A(t)=1.
\end{equation}
If the total number is $K$, the time to produce the first $k$ triples is
$\widetilde O(p+k)$ for every $1\le k\le K$, and the time to report
exhaustion is $\widetilde O(p+K)$. The algorithm is deterministic and uses
only polynomially many $O(\log p)$-bit words. The notation
$\widetilde O$ hides fixed powers of logarithms, uniformly over inputs.

This is an \textbf{unproved reporting hypothesis}. Ordinary fast Fourier
transform convolution does not supply it. In particular a total-output
bound without the prefix guarantee is not enough for the argument below.

\paragraph{Attempt: reduction theorem.}
\textbf{Conditional lemma.} If $\mathsf{StreamConv}$ holds, the problem in
this section has a uniform Las Vegas algorithm with worst-case expected
running time $\widetilde O(n^{3/2})$. Thus this hypothesis would disprove
the exact 3SUM conjecture printed above.

\emph{Construction.} For a valid input $S$ let
$m=\lceil n^{3/2}\rceil$. Enumerate the primes in $[m,2m]$ by a sieve and
choose one uniformly. The finitely many small $n$ are handled by brute
force. For all sufficiently large $m$ the prime number theorem gives
\begin{equation}\label{30006762:prime-count}
 \#\{p\text{ prime}:m\le p\le2m\}\gg m/\log m.
\end{equation}
The sieve and the explicit list use $\widetilde O(m)$ word operations and
space. An exactly uniform random index can be sampled by rejection from
a power-of-two range; the acceptance probability exceeds $1/2$.

For each residue $r$, form its bucket
$S_r=\{a\in S:a\equiv r\pmod p\}$ and set
$A(r)=1$ exactly for nonempty buckets. Invoke $\mathsf{StreamConv}$.
For each emitted triple $(r,s,t)$ inspect the Cartesian product
$S_r\times S_s\times S_t$, stopping immediately if
\[
 a+b+c=0\quad\text{as integers},\qquad |\{a,b,c\}|=3.
\]
If the reporter exhausts its list without finding this, answer no.
Every genuine zero triple has residue sum zero for every chosen prime,
so no yes instance can be missed. Every yes answer is checked in the
integers. The randomness affects only the cost, not correctness.

\paragraph{Attempt: a uniform expected-cost bound.}
Call an ordered integer triple \emph{false} if $a+b+c\ne0$ but
$p\mid a+b+c$. For such a triple set $q=a+b+c$; then
$0<|q|\le3n^3$. The number of different prime divisors of $q$ in
$[m,2m]$ is at most
\begin{equation}\label{30006762:divisor-count}
 \frac{\log(3n^3)}{\log m}\le3
\end{equation}
for all sufficiently large $n$. Indeed, the product of $k$ such distinct
primes is at least $m^k$ and divides $|q|$. Together with
\eqref{30006762:prime-count}, this gives
\begin{equation}\label{30006762:false}
 \E F_p
 \le \sum_{(a,b,c)\in S^3\, :\,a+b+c\ne0}
       \Pr(p\mid a+b+c)
 \ll \frac{n^3\log m}{m},
\end{equation}
where $F_p$ is the number of false ordered triples. No independence among
these events is used: this is only linearity of expectation.

Repeated-element zero triples need separate treatment. For each of the
three equality patterns, such as $a=b$, the equation $2a+c=0$ leaves at
most $n$ possibilities because $S$ is a set. Hence at most $3n$ ordered
zero triples have a repeated element. Counting the all-equal triple more
than once only strengthens this upper bound.

Before termination, therefore, at most
\begin{equation}\label{30006762:prefix-count}
 F_p+3n+1
\end{equation}
integer triples are inspected. On a no instance the last $1$ can be
omitted; on a yes instance there is only one inspected distinct-element
zero triple, because the algorithm stops at the first one. Each emitted
residue triple has a nonempty Cartesian product and accounts for at
least one inspected integer triple, except that the final product may
be only partly inspected. Thus the number of reporter outputs needed is
also at most \eqref{30006762:prefix-count}.

The prefix and exhaustion guarantees now imply cost
\[
 \widetilde O(m+n+F_p),\qquad
 \sup_{|S|=n}\E T(S)
 =\widetilde O\left(m+n+\frac{n^3}{m}\right)
 =\widetilde O(n^{3/2}).
\]
All integer sums, residues, array indices and sieve indices use
$O(\log n)$ bits. The constants in the estimates are independent of the
input values. Rejection sampling terminates almost surely, and all
remaining loops are finite. These observations check the uniform
Las Vegas and worst-case-expectation requirements, rather than replacing
them by a high-probability Monte Carlo statement.

\paragraph{Stress test: attempted implementation of the reporter.}
Convolution of $A$ with itself gives, for each $u$, the number of pairs
$r+s=u$. It does not give the first such pair for every requested target
in near-linear total time. Scanning $r$ independently for each occupied
$t$ can cost $p\,|\operatorname{supp}A|$; at
$p\asymp n^{3/2}$ this loses the claimed exponent. Building all residue
pairs explicitly costs $\Theta(|\operatorname{supp}A|^2)$ and also loses
the desired bound in the worst case. Hashing by itself therefore does
not finish the algorithm.

One might recursively bisect the support and use convolution to prune
empty pair rectangles. The unresolved amortization is that many target
sums can keep overlapping rectangles alive; I did not prove an
$\widetilde O(p+k)$ prefix bound on the total convolution work. Paying
for all $K$ outputs before emitting the first is invalid: a symmetric
set can have $\Theta(n^2)$ distinct-element zero triples. The quadratic
preprocessing in \eref{30006762}{2} is an explicit illustration of why a fast
query guarantee does not automatically solve the present problem.

The checking script exhausts all $462$ five-element subsets of
$[-5,5]$, using each of five small prime moduli. A deliberately slow
reporter confirms correctness and the handling of repeated entries.
This verifies neither $\mathsf{StreamConv}$ nor the asymptotic running
time. It also checks the elementary divisor bound at a finite sample
scale. No finite test could certify the required uniform reporting
complexity.

\paragraph{Dependency check.}
The reduction is not a lower bound against an unrestricted word-RAM and
is not a proof of the catalogue's conjecture. Its direction is the
opposite: the proposed reporting lemma would violate that conjecture,
with a polynomial exponent saving (for example any fixed saving below
$1/2$ after absorbing logarithms). Since 3SUM already belongs to
polynomial time, this consequence alone does not imply
$\mathrm P=\mathrm{NP}$ or its negation. The strong consequence is
nevertheless a warning against accepting the reporter as a routine FFT
primitive. The connection between integer 3SUM and output-sensitive
reporting is also consonant with the reductions in
\sref{30006762}{1}, Section 10.

\paragraph{Outcome.}
\textbf{Outcome: Conditional reduction.}

The hashing, false-positive estimate, repeated-element bound, stopping
argument, and exact model audit above form a proved conditional reduction.
The single missing algorithmic input is $\mathsf{StreamConv}$ with its
prefix-time and exhaustion guarantees. No implementation achieving those
guarantees has been obtained, and the finite checker is not such an
implementation. Thus no unconditional subquadratic algorithm or lower
bound is claimed.
% END RESEARCH ADDITION ATTEMPT-161
\subsection{Sources}
\begin{itemize}\raggedright
\item[\textnormal{[S1]}]\hypertarget{source:30006762:1}{} Abboud and Vassilevska Williams, Popular Conjectures Imply Strong Lower Bounds for Dynamic Problems (2014), Conjecture 1; KPP (2019) for Las Vegas convention.
\url{https://arxiv.org/pdf/1402.0054}.
{\small\textit{Catalogue formulation source.}}
% BEGIN RESEARCH ADDITION SOURCES-161
\item[\textnormal{[E1]}]\hypertarget{extra:30006762:1}{} Allan Gr\o nlund and Seth Pettie, ``Threesomes, Degenerates, and Love Triangles,'' arXiv:1404.0799 (2014). Introduction and 3SUM algorithm bounds.
\url{https://arxiv.org/abs/1404.0799}.
{\small\textit{Logarithmic improvements and the distinction from a fixed exponent saving.}}
\item[\textnormal{[E2]}]\hypertarget{extra:30006762:2}{} Shashwat Kasliwal, Adam Polak and Pratyush Sharma, ``3SUM in Preprocessed Universes: Faster and Simpler,'' \emph{TheoretiCS} 4 (2025), DOI 10.46298/theoretics.25.24; arXiv:2410.16784v3 (15 October 2025). Abstract, introduction and preprocessing/query theorem.
\url{https://arxiv.org/html/2410.16784v3}.
{\small\textit{Quadratic preprocessing is counted and is not available for free in the selected model.}}
% END RESEARCH ADDITION SOURCES-161
\end{itemize}

\end{document}
